% Part: proof-theory
% Chapter: natural-deduction
% Section: rules-N2

\documentclass[../../../include/open-logic-section]{subfiles}

\begin{document}

\begin{table}
\begin{tabular}{@{}c@{}c@{}}
\Axiom$x: !A, \Gamma \fCenter !B$
\RightLabel{\Intro{\lif}}
\UnaryInf$\Gamma \fCenter !A \lif !B $
\DisplayProof
&
\Axiom$\Gamma_1 \fCenter !A \lif !B$
\Axiom$\Gamma_2 \fCenter !A$
\RightLabel{\Elim{\lif}}
\BinaryInf$\Gamma_1, \Gamma_2 \fCenter !B$
\DisplayProof
\\[4ex]
\Axiom$\Gamma_1 \fCenter !A$
\Axiom$\Gamma_2 \fCenter !B$
\RightLabel{\Intro{\land}}
\BinaryInf$\Gamma_1, \Gamma_2 \fCenter !A \land !B $
\DisplayProof
&
\begin{tabular}[b]{@{}l@{}}
\Axiom$\Gamma \fCenter !A \land !B$
\RightLabel{\Elim{\land}[1]}
\UnaryInf$\Gamma \fCenter !A$
\DisplayProof
\\[3ex]
\Axiom$\Gamma \fCenter !A \land !B$
\RightLabel{\Elim{\land}[2]}
\UnaryInf$\Gamma \fCenter !B$
\DisplayProof\\[3ex]
\end{tabular}
\\[4ex]
\Axiom$\Gamma \fCenter !A$
\RightLabel{\Intro{\lor}[1]}
\UnaryInf$\Gamma \fCenter !A \lor !B$
\DisplayProof
&
\Axiom$ \Gamma \fCenter !B$
\RightLabel{\Intro{\lor}[2]}
\UnaryInf$\Gamma \fCenter !A \lor !B$
\DisplayProof\\[3ex]
\multicolumn{2}{c}{
\begin{tabular}[b]{@{}l@{}}
\Axiom$\Gamma_1 \fCenter !A \lor !B$
\Axiom$x:!A, \Gamma_2 \fCenter !C$
\Axiom$y:!B, \Gamma_3 \fCenter !C$
\RightLabel{\Elim{\lor}}
\TrinaryInf$\Gamma_1, \Gamma_2, \Gamma_3 \fCenter !C$
\DisplayProof
\end{tabular}
}
\\[4ex]
\Axiom$x:!A, \Gamma \fCenter \lfalse$
\RightLabel{\Intro{\lnot}}
\UnaryInf$ \Gamma \fCenter \lnot !A $
\DisplayProof
&
\Axiom$\Gamma_1 \fCenter \lnot !A $
\Axiom$\Gamma_2 \fCenter !A $
\RightLabel{\Elim{\lnot}}
\BinaryInf$\Gamma_1, \Gamma_2 \fCenter \lfalse$
\DisplayProof
\\[4ex]
\Axiom$\Gamma \fCenter !A(a) $
\RightLabel{\Intro{\lforall}}
\UnaryInf$\Gamma \fCenter \lforall[x][!A(x)]$
\DisplayProof
&
\Axiom$\Gamma \fCenter \lforall[x][!A(x)]$
\RightLabel{\Elim{\lforall}}
\UnaryInf$\Gamma \fCenter !A(t)$
\DisplayProof
\\[4ex]
\Axiom$\Gamma \fCenter !A(t) $
\RightLabel{\Intro{\lexists}}
\UnaryInf$\Gamma \fCenter \lexists[x][!A(x)]$
\DisplayProof
&
\Axiom$\Gamma_1 \fCenter  \lexists[x][!A(x)]$
\Axiom$x: !A(a), \Gamma_2 \fCenter !C$
\RightLabel{\Elim{\lexists}}
\BinaryInf$\Gamma_1, \Gamma_2 \fCenter !C$
\DisplayProof
\\[4ex]
\Axiom$\Gamma \fCenter \lfalse$
\RightLabel{\FalseInt}
\UnaryInf$\Gamma \fCenter !A$
\DisplayProof
&
\Axiom$x: \lnot !A, \Gamma \fCenter \lfalse$
\RightLabel{\FalseCl}
\UnaryInf$\Gamma \fCenter !A$
\DisplayProof
\end{tabular}
\caption{\Log{N2c} ও \Log{N2i}-র বিধি। $\Intro{\lforall}$ ও
$\Elim{\lexists}$-এ !!{constant}~$a$ সিদ্ধান্ত-সিকোয়েন্টে
ঘটে না। $\Gamma_i$ হলো লেবেলযুক্ত !!{formula}s~$x :
!A$-র সেট। \Intro{\lforall} ও \Elim{\lexists}-এর
আইগেনচল-শর্তে $a$ সিদ্ধান্ত-সিকোয়েন্টে থাকে না;
\Elim{\lexists}-এ প্রধান অস্তিত্বসূচক পূর্বধারণাতেও থাকে না।
\Log{N2i}-তে \FalseCl নেই।}
% BN-SRC-548 TARGET-CORRECTION-BEGIN
\par\smallskip\noindent\textnormal{\small\textbf{উৎস-সংশোধন (BN-SRC-548):} প্রথম ক্যাপশনে সাধারণ forall/exists নির্দেশের বদলে এই অধ্যায়ে সংজ্ঞায়িত lforall/lexists বিধি-চিহ্ন বসানো হয়েছে।}
% BN-SRC-548 TARGET-CORRECTION-END
% BN-SRC-549 TARGET-CORRECTION-BEGIN
\par\smallskip\noindent\textnormal{\small\textbf{উৎস-সংশোধন (BN-SRC-549):} N2 অস্তিত্ব-নিষ্কাশনে সিদ্ধান্ত-সিকোয়েন্টের পাশাপাশি আইগেনচল প্রধান অস্তিত্বসূচক পূর্বধারণাতেও না-থাকার শর্তটি স্পষ্ট করা হয়েছে।}
% BN-SRC-549 TARGET-CORRECTION-END
\ollabel{tab:N2}
\end{table}

\end{document}
